\subsection*{4.1 Notation and tools}
Before going into our results, we formally define permutation branching programs along with the notation we will use in the rest of this section. These programs will have a more specialized form than when we introduced them in Section2, which we subsequently justify. We assume basic familiarity with permutations, and we write $\sigma_1\sigma_2$ as a shorthand for $\sigma_2\circ\sigma_1$.
\paragraph{Definition17.} Let $n,m,s\in\mathbb N$ and let $f:\{0,1\}^n\to\{0,1\}$ be a function such that $f(0\ldots0)=0$. A permutation branching program is a sequence $P=\pi_1\ldots\pi_{\ell}$, where each $\pi_j$ is a pair $\langle i_j,\sigma_j\rangle$ where $i_j\in[0..n]$ and $\sigma_j$ is a permutation of $[m]$. We refer to each $\pi_j$ as an instruction of $P$. The width of $P$ is $m$ and the length of $P$ is $\ell$.
For any $\alpha\in\{0,1\}^n$ we define $P(\alpha)$ as follows: fix $\sigma=\mathrm{id}$, and for every $j=1\ldots s$, we set $\sigma$ to $\sigma\sigma_j$ if $\pi_j=\langle0,\sigma_j\rangle$ or $\pi_j=\langle i_j,\sigma_j\rangle$ where $\alpha_{i_j}=1$, and leave $\sigma$ unchanged otherwise (that is, if $\pi_j=\langle i_j,\sigma_j\rangle$ where $\alpha_{i_j}=0$). Our output is the final value of $\sigma$. We say that $P$ computes $f$ if there exists a permutation $\sigma^*\neq\mathrm{id}$ such that $P(\alpha)=\mathrm{id}$ if $f(\alpha)=0$ and $P(\alpha)=\sigma^*$ if $f(\alpha)=1$.
We make three observations about our choices in Definition17. First, the restriction that $f(0\ldots0)=0$ will be a convenience; we can always compute $\neg f$ instead if this condition does not hold, or change Definition17 such that $P(\alpha)=\mathrm{id}$ if $f(\alpha)=1$ and vice versa.\footnote{We also note that if $P$ computes $\neg f$, we can compute $f$ by appending the instruction $\langle0,(\sigma^*)^{-1}\rangle$ to $P$. We avoid taking this route because a later observation will allow us to remove these fixed layers, but only when $f(\alpha)=0$, which would cause our logic to become circular.} Second, in a layer $\langle i,\sigma_j\rangle$ reading variable $x_i$, we only fix a permutation in the case that $x_i$ is set to1. This is without loss of generality, as adding a layer of the form $\langle0,\sigma'_j\rangle$ before an instruction can be thought of as choosing a permutation in the case that $x_i$ is set to0 (while the permutation for $x_i=1$ can be adjusted accordingly).
Before going on to our third observation, we state and prove four simple lemmas which will allow us to conveniently restructure our programs $P$.
\paragraph{Lemma18.} Let $P$ be a permutation branching program computing $f$ and let $j$ be such that $i_j=i_{j+1}$. Then the program $P'$ resulting from replacing $\pi_j,\pi_{j+1}$ with $\pi'_j=\langle i_j,\sigma_j\sigma_{j+1}\rangle$ is also a valid program for computing $f$.
\begin{proof}
In both $P$ and $P'$, the permutations $\sigma_j$ and $\sigma_{j+1}$ are both applied when $i_j=1$ and neither are applied when $i_j=0$.
\end{proof}
\paragraph{Lemma19.} Let $P$ be a permutation branching program computing $f$ and let $j$ be such that $\sigma_j\sigma_{j+1}=\sigma_{j+1}\sigma_j$. Then the program $P'$ resulting from switching the order of $\pi_j$ and $\pi_{j+1}$ is also a valid program for computing $f$.
\begin{proof}
Consider any assignment $\alpha$ to $\mathbf x$. In the case that either $\alpha_{i_j}$ or $\alpha_{i_{j+1}}$ is set to0, these programs compute identical permutations as either $\sigma_j$ or $\sigma_{j+1}$ will not be applied. If both are set to1, then
\[P'(\alpha)=\Sigma_1\sigma_{j+1}\sigma_j\Sigma_2=\Sigma_1\sigma_j\sigma_{j+1}\Sigma_2=P(\alpha)\]
where $\Sigma_1,\Sigma_2$ are the permutations corresponding to the rest of the instructions on input $\alpha$.
\end{proof}
\paragraph{Lemma20.} Let $P=\pi_1\ldots\pi_s$ be a permutation branching program computing $f$, let $\pi_j=\langle i_j,\sigma_j\rangle$ for all $j$, and let $j^*\in[\ell]$ be such that $i_{j^*}=0$. Then there exists a permutation branching program $P'=\pi'_1\ldots\pi'_{j^*-1}\pi'_{j^*+1}\ldots\pi'_{\ell}\pi_{j^*}$ computing $f$, where $\pi'_j=\langle i_j,\sigma'_j\rangle$ for some permutation $\sigma'_j$.
\begin{proof}
For $j<j^*$ define $\sigma'_j=\sigma_j$, and for $j>j^*$ define $\sigma'_j=\sigma_{j^*}\sigma_j\sigma_{j^*}^{-1}$. Clearly because $\sigma_{j^*}^{-1}$ and $\sigma_{j^*}$ cancel out for every adjacent pair of permutations $\sigma'_j$, $P'(\alpha)$ contains exactly the same permutations as $P(\alpha)$ in the same order regardless of the assignment $\alpha$.\footnote{This argument actually allows us to move $\pi_{j^*}$ to any spot in the program we want, but we are content with just moving them to the end, for reasons which will become immediately clear.}
\end{proof}
Our next observation is that the layers of the form $\langle0,\sigma_j\rangle$ are only a convenience and are not necessary to our definition. Let $P$ be our program for $f$ and let $\sigma_1\ldots\sigma_k$ be the permutations corresponding to instructions $\pi_j$ where $i_j=0$. By our restriction that $f(0\ldots0)=0$ we get $\sigma_1\ldots\sigma_k=P(0\ldots0)=\mathrm{id}$, and by Lemma20 we can move the instructions $\pi_j$ with $i_j=0$ to the end of the program, in order, at which point we can simply remove them all using Lemma18 as they compose to the identity for any input $\alpha$.
We can also generalize Lemma20 for restrictions of the function $f$, meaning when we fix the values of some variables and consider the function on the remaining variables. This is simply the observation that fixing variable $x_i$ turns all instructions of the form $\langle i,\sigma_j\rangle$ into fixed layers $\langle0,\sigma_j\rangle$.
\paragraph{Corollary21.} Let $\rho\in\{0,1,*\}^n$ and let $f_{\rho}$ be the function $f$ with $x_i$ fixed to $\rho(i)$ wherever $\rho(i)\neq*$. Let $P=\pi_1\ldots\pi_{\ell}$ be a permutation branching program computing $f$, let $\pi_j=\langle i_j,\sigma_j\rangle$ for all $j$, and let $j^*\in[s]$ be such that $i_{j^*}=0$ or $\rho(i_{j^*})\neq*$. Then there exists a permutation branching program $P'=\pi'_1\ldots\pi'_{j^*-1}\pi'_{j^*+1}\ldots\pi'_{\ell}\pi_{j^*}$ computing $f_{\rho}$, where $\pi'_j=\langle i'_j,\sigma'_j\rangle$ for some permutation $\sigma'_j$ and $i'_j=i_j$ iff $\rho(i_j)=*$ and0 otherwise.
\begin{proof}
Let program $P''$ be the result of replacing $i_j$ with0 in each instruction $\pi_j\in P$ such that $\rho(i_j)\neq*$. Clearly this program computes $f_{\rho}$, and so applying Lemma20 to $P''$ completes the proof.
\end{proof}
Assuming that $f_{\rho}(0\ldots0)=0$, by our previous observation this allows us to remove all layers that read variables fixed by $\rho$. We also note that the other three lemmas hold for $f_{\rho}$ with no changes.
Finally, when $f$ is the AND function, we can rotate our program by moving the first instruction to the end as many times as we like:
\paragraph{Lemma22.} Let $P=\Pi_1,\Pi_2$ be a permutation branching program computing $\mathrm{AND}(x_1,\ldots,x_n)$, where $\Pi_1$ and $\Pi_2$ are sequences of instructions. Then $P'=\Pi_2,\Pi_1$ also computes $\mathrm{AND}(x_1,\ldots,x_n)$.
\begin{proof}
Let $\sigma^*=P'(1,\ldots,1)$ (that is, $P'$ applied to the all-ones input). Referring to Definition17, to prove that $P'$ computes AND, we must show(a) that $\sigma^*\neq\mathrm{id}$,(b) that $P'(\alpha)=\mathrm{id}$ whenever $\mathrm{AND}(\alpha)=0$, and(c) that $P'(\alpha)=\sigma^*$ whenever $\mathrm{AND}(\alpha)=1$.
(c) follows from the fact that the only string $\alpha$ satisfying $\mathrm{AND}(\alpha)=1$ is the all-ones string.
Now, let $\alpha\in\{0,1\}^n$. Note that $P(\alpha)=\Pi_1(\alpha)\Pi_2(\alpha)$. If $\mathrm{AND}(\alpha)=0$, then $P(\alpha)=\mathrm{id}$, so $\Pi_1(\alpha)=(\Pi_2(\alpha))^{-1}$, and so $P'(\alpha)=\mathrm{id}$ as well: this proves(b).
Finally, note that $P(1,\ldots,1)\neq\mathrm{id}$, so $\Pi_1(1,\ldots,1)\neq(\Pi_2(1,\ldots,1))^{-1}$, and so $\sigma^*=P'(1,\ldots,1)\neq\mathrm{id}$, proving(a).
\end{proof}
